\chapter{Allergic Reactions to Cow's Milk}

\section*{Definition and Overview}
A cow's milk allergy is an adverse immune-mediated reaction to proteins found in cow's milk, chiefly casein and the whey proteins alpha-lactalbumin and beta-lactoglobulin. It is one of the most common food allergies of infancy and early childhood, and it must be distinguished clearly from lactose intolerance, which is a digestive enzyme deficiency rather than an immune disorder. Cow's milk allergy may present in two broad patterns: an immediate (IgE-mediated) form in which symptoms appear within minutes to a couple of hours of exposure, and a delayed (non-IgE-mediated) form in which symptoms develop hours to days later, frequently involving the gastrointestinal tract. The cornerstone of management is strict avoidance of cow's milk protein in all its forms, together with nutritional support to ensure normal growth and development. Because the allergy is commonly outgrown, regular specialist review to reassess tolerance is an important part of long-term care.

\section*{Epidemiology}
Cow's milk allergy is reported in approximately 1 to 3\% of infants in developed countries, making it among the most frequent food allergies in early life. It typically presents in the first year, often within weeks to months of the introduction of cow's milk formula. The natural history is favourable in most children: the non-IgE-mediated form frequently resolves by three to five years of age, whereas the IgE-mediated form persists longer and may continue into adolescence or adulthood in a significant minority. Some population-based studies suggest that the prevalence of parent-reported allergy is considerably higher than the prevalence confirmed by formal challenge testing, highlighting the importance of accurate diagnosis. Cow's milk allergy in adults is uncommon and usually represents persistence of childhood disease. Risk factors include a personal or family history of atopy, other food allergies, eczema, and coexisting asthma or allergic rhinitis.

\section*{Aetiology and Pathophysiology}
The underlying abnormality is an inappropriate immune response directed against cow's milk proteins, most commonly beta-lactoglobulin, alpha-lactalbumin, and casein. Two principal mechanisms are recognised:

\begin{itemize}
    \item \textbf{IgE-mediated allergy:} allergen-specific immunoglobulin E is produced and binds to mast cells and basophils; on re-exposure, the milk proteins cross-link these antibodies and trigger release of histamine and other mediators, producing symptoms within minutes to two hours.
    \item \textbf{Non-IgE-mediated allergy:} T-cell-driven inflammatory responses within the gut lead to delayed symptoms such as vomiting, reflux, diarrhoea, blood in the stool, constipation, and eczema, typically appearing hours to days after ingestion.
    \item \textbf{Mixed and indeterminate patterns:} some children have features of both mechanisms, and the clinical expression is influenced by the quantity consumed, the state of the intestinal barrier, and coexisting sensitisation to other foods.
\end{itemize}

The diagnosis of the non-IgE form is particularly difficult because there is no reliable laboratory marker, and it rests on clinical response to dietary elimination and rechallenge.

\section*{Clinical Features}
Presentation depends on the immune mechanism, the age of the child, and the amount of milk protein ingested:

\begin{itemize}
    \item \textbf{Immediate reactions:} urticaria, erythema, itching, swelling of the lips, tongue, or face, vomiting, acute rhinitis, and wheezing shortly after ingestion.
    \item \textbf{Delayed reactions:} recurrent regurgitation and reflux-like behaviour, loose or blood-streaked stools, constipation, infantile colic, eczema flares, and poor weight gain.
    \item \textbf{Severe reactions:} laryngeal oedema with stridor or a hoarse voice, bronchospasm, hypotension, and collapse, representing anaphylaxis.
    \item \textbf{Non-specific features:} irritability, poor sleep, and food refusal, which are common but non-specific in infants.
\end{itemize}

The threshold dose required to trigger a reaction varies greatly between individuals, and even small amounts of hidden milk protein in processed foods can provoke symptoms in highly sensitive children.

\section*{Differential Diagnosis}
Several conditions may mimic cow's milk allergy:

\begin{itemize}
    \item \textbf{Lactose intolerance:} bloating, wind, and diarrhoea after milk due to lactase deficiency; not an immune reaction and does not cause skin or respiratory symptoms.
    \item \textbf{Other food allergies:} reactions to soy, egg, wheat, or fish may coexist with or be mistaken for milk allergy.
    \item \textbf{Gastro-oesophageal reflux disease:} a common cause of regurgitation in infants, unrelated to allergy.
    \item \textbf{Gastroenteritis and gut infections:} may cause transient vomiting and diarrhoea.
    \item \textbf{Coeliac disease and malabsorption:} cause poor growth and chronic diarrhoea.
    \item \textbf{Non-disease causes:} teething, nappy rash, and normal infant fussiness are frequently misinterpreted as allergy by parents.
\end{itemize}

\section*{When to Seek Medical Care}
Parents should seek medical review if a baby or child has persistent vomiting, diarrhoea, blood in the stool, poor weight gain, troublesome eczema, or symptoms that repeatedly follow feeds. Earlier review is advisable when symptoms begin soon after the first introduction of cow's milk formula.

\textbf{Call 000 (or local emergency services) for:}
\begin{itemize}
    \item Difficulty breathing, wheezing, or a hoarse voice after drinking milk.
    \item Swelling of the face, tongue, or throat, or difficulty swallowing.
    \item Collapse, floppiness, or loss of consciousness.
    \item Repeated forceful vomiting with signs of dehydration, such as a dry mouth, sunken eyes, or reduced urine output.
\end{itemize}

\section*{Diagnosis and Investigations}
Diagnosis relies primarily on a detailed clinical history supported by targeted tests:

\begin{itemize}
    \item \textbf{History:} the timing and nature of symptoms relative to milk ingestion, the type of formula used, and any family history of atopy.
    \item \textbf{Elimination diet:} strict removal of cow's milk protein from the diet for two to four weeks, with observation of symptom improvement.
    \item \textbf{Oral food challenge:} supervised reintroduction of milk to confirm the diagnosis; challenges for IgE-mediated disease should be performed in a facility equipped to manage anaphylaxis.
    \item \textbf{Skin prick testing and specific IgE:} useful for immediate-type reactions but not diagnostic on their own, because sensitisation does not always equal clinical allergy.
    \item \textbf{Atopy patch testing:} used in some specialist centres for delayed gastrointestinal reactions, though its routine value is debated.
    \item \textbf{Referral:} to a paediatrician or allergy specialist is appropriate for complex, severe, or diagnostic-uncertain cases, and for planning formal challenges.
\end{itemize}

\section*{Management}
Management centres on the complete removal of cow's milk protein from the diet while maintaining adequate nutrition:

\begin{itemize}
    \item \textbf{Breastfed infants:} the mother may need to avoid cow's milk in her own diet, under medical and dietetic guidance, with attention to her calcium and vitamin D intake.
    \item \textbf{Bottle-fed infants:} an extensively hydrolysed formula is first-line; amino-acid formulas are reserved for infants who fail extensively hydrolysed formula or who have severe reactions.
    \item \textbf{Older children and adults:} strict avoidance of milk, cheese, butter, yoghurt, and cream, together with careful reading of ingredient labels for casein, whey, milk solids, and lactose-free products that still contain milk protein.
    \item \textbf{Anaphylaxis action plan:} a written plan covering avoidance, recognition of severe symptoms, and use of an adrenaline auto-injector when prescribed.
    \item \textbf{Dietetic review:} regular review to ensure adequate calcium, protein, energy, and vitamin D for growth, and to prevent nutritional deficiencies.
    \item \textbf{Education:} informing childcare, school, restaurants, and family members, and teaching the child age-appropriate self-management skills as they grow older.
\end{itemize}

\section*{Complications}
Potential complications of cow's milk allergy and its management include:

\begin{itemize}
    \item Poor weight gain and faltering growth if the diet is not adequately planned.
    \item Accidental exposure causing acute reactions, occasionally severe or life-threatening.
    \item Nutritional deficiencies, particularly of calcium and vitamin D, when elimination is unsupervised.
    \item Chronic eczema, disturbed sleep, and feeding difficulties in affected infants.
    \item Anxiety about feeding in parents and children, and social restrictions imposed by the diet.
\end{itemize}

\section*{Prognosis and Outcomes}
The prognosis is generally excellent. Most children with non-IgE-mediated cow's milk allergy outgrow it by three to five years of age, and a substantial proportion of children with the IgE-mediated form also develop tolerance by school age, although persistence into adolescence and adulthood is more likely in this group. Higher specific IgE levels, coexisting asthma, and other food allergies are associated with slower resolution. Children who outgrow the allergy can usually reintroduce milk into the diet, ideally under specialist supervision, and with careful avoidance and nutritional support growth and development are normally unaffected.

\section*{Prevention and Self-Care}
There is no reliably proven strategy for preventing the first development of cow's milk allergy, but several measures may help:

\begin{itemize}
    \item \textbf{Breastfeeding} where possible, as it supports normal infant immune development and overall health.
    \item \textbf{Introduction of foods} at the recommended age, including cow's milk products in appropriate forms, under medical advice.
    \item \textbf{Reading food labels} carefully and learning the many names under which milk protein appears.
    \item \textbf{Communicating the allergy} to childcare providers, schools, and family members.
    \item \textbf{Carrying prescribed medicines}, including an adrenaline auto-injector, together with an up-to-date action plan.
    \item \textbf{Regular follow-up} with an allergy specialist to reassess tolerance and adjust management as the child grows.
\end{itemize}

\section*{Sources and Further Reading}
The following reputable sources provide reliable, up-to-date information on cow's milk allergy:

\begin{itemize}
    \item NHS. \textit{Cow's milk allergy information for parents and healthcare professionals.} https://www.nhs.uk/
    \item Mayo Clinic. \textit{Milk allergy: symptoms, causes, and treatment.} https://www.mayoclinic.org/
    \item MSD Manual. \textit{Food allergy overview, including cow's milk allergy.} https://www.msdmanuals.com/
    \item MedlinePlus. \textit{Milk allergy health information.} https://medlineplus.gov/
    \item NICE. \textit{Clinical guidance on food allergy in children and young people.} https://www.nice.org.uk/
    \item World Health Organization. \textit{Information on breastfeeding and infant nutrition.} https://www.who.int/
\end{itemize}
